\documentclass[12pt,a4paper]{article}

\usepackage[utf8]{inputenc}
\usepackage[T1]{fontenc}
\usepackage[spanish,es-nodecimaldot]{babel}
\usepackage{geometry}
\usepackage{xcolor}
\usepackage{listings}
\usepackage{parskip}
\usepackage{needspace}

\geometry{margin=2.5cm}

\definecolor{codigo}{RGB}{245,245,245}
\definecolor{azul}{RGB}{30,70,130}

\lstset{
    language=Java,
    backgroundcolor=\color{codigo},
    basicstyle=\ttfamily\small,
    keywordstyle=\color{azul}\bfseries,
    commentstyle=\color{gray},
    showstringspaces=false,
    frame=single,
    breaklines=true,
    columns=fullflexible
}

\title{Informe sobre los principios SOLID}
\author{Jorman Noriega}
\date{\today}

\begin{document}

\maketitle

\begin{abstract}
Este informe explica de forma sencilla los cinco principios SOLID de la
programación orientada a objetos. Los ejemplos utilizan situaciones de una
tienda para mostrar cómo escribir código más organizado y fácil de mantener.
\end{abstract}

\section{Introducción}

SOLID es un conjunto de cinco principios para diseñar programas orientados a
objetos. Su objetivo es reducir la dependencia entre las clases, facilitar los
cambios y evitar que el código se vuelva difícil de mantener.

Las letras de SOLID representan los siguientes principios:

\begin{itemize}
    \item S: Single Responsibility Principle.
    \item O: Open/Closed Principle.
    \item L: Liskov Substitution Principle.
    \item I: Interface Segregation Principle.
    \item D: Dependency Inversion Principle.
\end{itemize}

\section{S: Responsabilidad única}

Una clase debe tener una sola responsabilidad y una sola razón para cambiar.
Por ejemplo, una clase \texttt{Pedido} debería encargarse de los datos del
pedido, pero no de calcular totales y guardarlos en una base de datos al mismo
tiempo.

\textbf{Ejemplo incorrecto:}

\Needspace{12\baselineskip}
\begin{lstlisting}
class Pedido {
    void agregarProducto() { }
    void calcularTotal() { }
    void guardarEnBaseDatos() { }
}
\end{lstlisting}

La clase anterior mezcla tres responsabilidades diferentes.

\textbf{Ejemplo correcto:}

Una solución sencilla es separar las responsabilidades:

\Needspace{18\baselineskip}
\begin{lstlisting}
class Pedido {
    private double total;
    public double getTotal() { return total; }
}

class CalculadoraPedido {
    public double calcular(Pedido pedido) {
        return pedido.getTotal();
    }
}

class PedidoRepository {
    public void guardar(Pedido pedido) {
        System.out.println("Pedido guardado");
    }
}

public class Main {
    public static void main(String[] args) {
        Pedido pedido = new Pedido();
        CalculadoraPedido calculadora = new CalculadoraPedido();
        System.out.println(calculadora.calcular(pedido));
        new PedidoRepository().guardar(pedido);
    }
}
\end{lstlisting}

De esta forma, cada clase tiene una función clara. Si cambia la forma de
guardar los pedidos, no es necesario modificar la clase \texttt{Pedido}.

\section{O: Abierto/cerrado}

Una clase debe estar abierta para extenderse, pero cerrada para modificarse.
Esto significa que podemos agregar comportamientos nuevos sin cambiar el
código que ya funciona.

\textbf{Ejemplo incorrecto:}

\Needspace{12\baselineskip}
\begin{lstlisting}
void pagar(String tipo, double valor) {
    if (tipo.equals("tarjeta")) { }
    else if (tipo.equals("paypal")) { }
    // Cada nuevo pago modifica esta funcion.
}
\end{lstlisting}

\textbf{Ejemplo correcto:} se pueden agregar diferentes formas de pago
mediante una interfaz:

\Needspace{18\baselineskip}
\begin{lstlisting}
interface MetodoPago {
    void pagar(double valor);
}

class TarjetaPago implements MetodoPago {
    public void pagar(double valor) {
        System.out.println("Pago con tarjeta");
    }
}

class NequiPago implements MetodoPago {
    public void pagar(double valor) {
        System.out.println("Pago con Nequi");
    }
}

public class Main {
    public static void main(String[] args) {
        MetodoPago pago = new NequiPago();
        pago.pagar(25000);
    }
}
\end{lstlisting}

Para agregar otro método, como PayPal, se crea una nueva clase sin modificar
las clases anteriores.

\section{L: Sustitución de Liskov}

Una clase hija debe poder utilizarse en cualquier lugar donde se espere su
clase padre sin causar errores ni comportamientos inesperados.

\textbf{Ejemplo incorrecto:}

\Needspace{16\baselineskip}
\begin{lstlisting}
class Empleado {
    public double calcularSalario() { return 2000000; }
}

class EmpleadoMedioTiempo extends Empleado {
    public double calcularSalario() {
        throw new UnsupportedOperationException();
    }
}

// Esta llamada provoca un error:
mostrarSalario(new EmpleadoMedioTiempo());
\end{lstlisting}

El empleado de medio tiempo no puede sustituir a \texttt{Empleado}, porque
rompe el comportamiento esperado. Esto viola LSP.

\textbf{Ejemplo correcto:} los dos tipos de empleado pueden utilizarse como un
\texttt{Empleado}:

\Needspace{18\baselineskip}
\begin{lstlisting}
abstract class Empleado {
    public abstract double calcularSalario();
}

class EmpleadoTiempoCompleto extends Empleado {
    public double calcularSalario() { return 2000000; }
}

class EmpleadoMedioTiempo extends Empleado {
    public double calcularSalario() { return 1000000; }
}

public class Main {
    public static void main(String[] args) {
        mostrarSalario(new EmpleadoTiempoCompleto());
        mostrarSalario(new EmpleadoMedioTiempo());
    }

    static void mostrarSalario(Empleado empleado) {
        System.out.println(empleado.calcularSalario());
    }
}
\end{lstlisting}

Una función que reciba un \texttt{Empleado} puede trabajar con cualquiera de
las dos clases hijas sin conocer sus detalles.

\section{I: Segregación de interfaces}

Una clase no debe estar obligada a implementar métodos que no necesita. Para
esto, es mejor tener varias interfaces pequeñas en lugar de una interfaz muy
grande.

\textbf{Ejemplo incorrecto:} una interfaz grande obliga al robot a implementar
métodos que no puede usar.

\Needspace{16\baselineskip}
\begin{lstlisting}
interface Trabajador {
    void trabajar();
    void comer();
    void descansar();
}

class Robot implements Trabajador {
    public void trabajar() { }
    public void comer() { /* No aplica */ }
    public void descansar() { /* No aplica */ }
}
\end{lstlisting}

\textbf{Ejemplo correcto:} se divide la interfaz grande en interfaces pequeñas.

\Needspace{18\baselineskip}
\begin{lstlisting}
interface Trabajable { void trabajar(); }
interface Comible { void comer(); }
interface Descansable { void descansar(); }

class Empleado implements Trabajable, Comible, Descansable {
    public void trabajar() { System.out.println("Empleado trabaja"); }
    public void comer() { System.out.println("Empleado come"); }
    public void descansar() { System.out.println("Empleado descansa"); }
}

class Robot implements Trabajable {
    public void trabajar() { System.out.println("Robot trabaja"); }
}

public class Main {
    public static void main(String[] args) {
        new Empleado().comer();
        new Robot().trabajar();
    }
}
\end{lstlisting}

El robot solo implementa \texttt{Trabajable}; no depende de \texttt{Comible} ni
de \texttt{Descansable}.

\section{D: Inversión de dependencias}

Las clases principales deben depender de abstracciones, como interfaces, y no
directamente de clases concretas.

\textbf{Ejemplo incorrecto:} el servicio depende directamente de una clase
concreta.

\Needspace{12\baselineskip}
\begin{lstlisting}
class PedidoService {
    void confirmarPedido() {
        new EmailNotificador().enviar("Pedido confirmado");
    }
}
\end{lstlisting}

\textbf{Ejemplo correcto:} el servicio puede recibir cualquier tipo de
notificador:

\Needspace{18\baselineskip}
\begin{lstlisting}
interface Notificador {
    void enviar(String mensaje);
}

class EmailNotificador implements Notificador {
    public void enviar(String mensaje) {
        System.out.println("Correo: " + mensaje);
    }
}

class PedidoService {
    private Notificador notificador;

    public PedidoService(Notificador notificador) {
        this.notificador = notificador;
    }

    public void confirmarPedido() {
        notificador.enviar("Pedido confirmado");
    }
}

public class Main {
    public static void main(String[] args) {
        Notificador notificador = new EmailNotificador();
        PedidoService servicio = new PedidoService(notificador);
        servicio.confirmarPedido();
    }
}
\end{lstlisting}

Así, \texttt{PedidoService} puede trabajar con un notificador por correo,
WhatsApp u otro medio sin modificar su código.

\section{Conclusión}

Los principios SOLID ayudan a construir programas más ordenados, flexibles y
fáciles de probar. No es necesario aplicar soluciones complejas: separar
responsabilidades, usar interfaces y evitar dependencias directas ya mejora
considerablemente la calidad del código.

En el proyecto de ejemplo, los cinco principios se aplican a diferentes
partes de una tienda: pedidos, pagos, empleados, impresoras y notificaciones.

\end{document}
